%%=====================================================================
%% M2 — WP4: the warp exponent and the value of ζ. SURGICAL FRAGMENT.
%% UK English. "We". Reproduces, in-document and without citation, the
%% derivation of Paper 2 §§3–4 and Paper I Addendum §§3–4 (fossil), on the
%% footing of M2 §3 (the base system is the restriction of the Master
%% Equation). PI rulings: Q12/Q14 (weak (c)), Q15 (direction: τ-symmetric,
%% j/k convention; τ₂ physical), Q16 (jump exact, constant slope a boundary
%% condition), Λ₁₂ by stationarity (M2-6). Placeholders: §3 = WP3, §2 = WP0/WP2.
%%=====================================================================

\section{The warp exponent and the value of $\zeta$}\label{sec:warp}

The purpose of this section is to solve the base restriction of the Master
Equation, \eqref{eq:2D}, for the profile of the warp across a half-period
of the pillowcase, and to show that the exponent by which the warp falls is
$N\pi$ with $N$ the number of flux units, that $N=\chi(\Keight)=12$, and
hence that the coupling of Section~\ref{sec:equation} has the value
$\zeta=12\pi$. We announce the two results.

\begin{theorem}[exact warp exponent per flux unit]\label{thm:warp-exponent}
For the base system \eqref{eq:2D} on $\Ttwo/\Ztwo$ with $N$ units of
quantised flux localised at a cone point, the warp exponent across the
half-period transverse to the wall on which the flux sits satisfies
exactly
\begin{equation}\label{eq:kL-Npi}
  kL\;=\;N\pi\,,
\end{equation}
where $k$ is the constant decay rate of the warp exponent between the
walls and $L$ the corresponding half-period. The factor $\pi$ is the
deficit angle at the cone point.
\end{theorem}

\begin{corollary}[the value of $\zeta$]\label{cor:zeta}
With $N=\chi(\Keight)=12$ (Proposition~\ref{prop:flux}) and the dilaton
normalised to fall by one unit across the half-period (\eqref{eq:a-Phi}),
\begin{equation}\label{eq:zeta}
  \zeta\;=\;kL\;=\;\chi(\Keight)\,\pi\;=\;12\pi\,.
\end{equation}
\end{corollary}

\subsection{The topological integers}\label{sec:integers}

\begin{proposition}\label{prop:chi-pillowcase}
$\chi(\Ttwo/\Ztwo)=2$, and the total curvature of the pillowcase is $4\pi$,
concentrated as a deficit of exactly $\pi$ at each of the four cone points.
\end{proposition}

\begin{proof}
We give two independent routes to the Euler characteristic, and then the
curvature. First route, by the cone points. The Euler characteristic of a
torus is zero. Folding it in half by the map $z\mapsto-z$ would halve that
to zero again, except at the four points the fold leaves fixed, where the
quotient is not smooth: each fixed point contributes a correction of
$1-\tfrac12=\tfrac12$, because a neighbourhood of it is a cone of angle
$\pi$, half a full turn. Four corrections of a half make two, so
$\chi=0+4\cdot\tfrac12=2$ (M1 Theorem~5.7, \texttt{L2\_chi\_pillowcase}).
Second route, by counting holes. The folded torus is, as a topological
space, a sphere --- stretch the pillowcase out and it is a ball with four
puckered points --- and a sphere has one connected piece, no independent
loops and one enclosed cavity, Betti numbers $(1,0,1)$, so $\chi=1-0+1=2$.
The two routes agree. For the curvature, the Gauss--Bonnet theorem for
orbifolds \cite{Kawasaki1979} says the total curvature is $2\pi$ times the
Euler characteristic, $4\pi$. The pillowcase is flat everywhere except at
the four cone points, so all $4\pi$ of curvature sits at them, $\pi$ at
each --- which is exactly the angle missing at a cone of angle $\pi$
(\texttt{L5\_deficit}). \rigR
\end{proof}

\begin{proposition}\label{prop:chi-K8}
$\chi(\Keight)=\chi(\CPtwo)\,\chi(\Stwo)\,\chi(\Ttwo/\Ztwo)=3\cdot2\cdot2=12$.
\end{proposition}

\begin{proof}
Multiplicativity of $\chi$ over fibrations with compact fibre; the warp
does not change the additive structure of cohomology
(M1 Theorem~7.6, \texttt{L3\_chi\_K8}). \rigR
\end{proof}

\begin{proposition}[the flux number]\label{prop:flux}
Tadpole cancellation on $\Keight$ fixes the number of quantised flux units
of the determinant line to $N_{\mathrm{flux}}=\chi(\Keight)=12$, and all
twelve sit on the base cycle.
\end{proposition}

\begin{proof}
On a compact space the total charge sourced must cancel. The gravitational
contribution to the tadpole is the Euler density, whose integral is
$\chi(\Keight)$; the flux contribution is integrally quantised and must
balance it, so $N_{\mathrm{flux}}=\chi(\Keight)=12$. Cross-check from the
parity structure: the individual quanta are odd on the non-spin factor
$\CPtwo$ and even on the spin factor $\Stwo$, by $c_{1}\equiv w_{2}\bmod2$
(M1 Proposition~7.3, \texttt{L4\_flux\_parity}), and the values
$(Q_{c},Q_{w})=(3,6)$ of M1 Proposition~8.2 are the ones the constraint
admits. The twelve units sit on the base because a localised flux
satisfying the first-order equation of Section~\ref{sec:bps} needs a cone
point to sit at, and the base is the only cycle that has them. The
smallest case: one unit at one corner balanced by one unit of integrated
Euler density on its cone neighbourhood; the general case is the sum of
twelve such balances. \rigD
\end{proof}
%% Weak (c), Q14: flux quanta are topological data; the placement argument is
%% the Addendum's §4.6 reason, reproduced. Minimality Register M2-2.

\subsection{The base system and its single modulus}\label{sec:base-system}

Section~\ref{sec:base-restriction} obtained the base restriction
\eqref{eq:2D} of the Master Equation: two-dimensional gravity with
coefficient $\kappa_{2}^{-2}$, the energy $\tfrac12\lvert F\rvert^{2}$ of
the line's field strength, a potential, and the gradient terms of the
fibre warp $a$ and the dilaton $\QCD$. Two facts reduce it to a system with
a single modulus.

\paragraph{The fibre decouples from the jump.}
$\CPtwo$ with the Fubini--Study metric and $\Stwo$ with the round metric
are Einstein, $\mathrm{Ric}=(6/R_{c}^{2})g$ and $\mathrm{Ric}=(1/R_{w}^{2})g$
(M1 Propositions~7.1, 7.2). For a warp depending only on the base
coordinates the warped-product curvature \cite{BishopONeill1969} therefore
contributes to the base equation only smooth scalar terms in $a$ --- a
potential and a modification of the gradient term --- with no mixed
base--fibre components. Smooth terms cannot change a jump condition at a
cone point; they change the shape of the profile between cone points.
The effective two-dimensional gravitational coefficient is the
fibre-volume-weighted one,
\begin{equation}\label{eq:kappa-eff}
  \kappa_{2}^{2}\;=\;\frac{\kappa_{12}^{2}}{V_{4}\,\mathrm{Vol}_{0}\,e^{6a_{0}}}\,,
  \qquad \mathrm{Vol}_{0}=\mathrm{Vol}(\CPtwo)\,\mathrm{Vol}(\Stwo)\,,
\end{equation}
with $V_{4}$ the external volume and $a_{0}$ the warp at the wall. \rigD

\paragraph{The dilaton and the warp are one modulus.}
On the solutions constructed below the dilaton's profile and the fibre
warp are tied by \eqref{eq:a-Phi}, $a=-\zeta\QCD$, with $\QCD$ the
canonically normalised modulus of the warped direction and $\QCD=1$ the
configuration at the far wall. This is the definition of the
normalisation of $\QCD$ on the base, and it is what makes the coupling
$\zeta$ of \eqref{eq:master} and the exponent of \eqref{eq:kL-Npi} one
number: a profile of $\QCD$ that falls by one unit across the half-period
is a warp exponent that falls by $\zeta$. In the fossil form of this
construction the transverse base length itself was written with the
factor $e^{2\zeta\QCD}$ while the fibre carried $e^{2a}$; under
\eqref{eq:a-Phi} these are two expressions of the same exponent and not
two warps. \rigD
%% Q16 reconciliation: Paper 2 §7.3 (a = −ζΦ, Φ = 1 boundary) adopted as the
%% normalisation; Addendum eq. 8D-metric's e^{2ζΦ}dθ_ω² read as the same
%% exponent on the transverse length. Flag: "canonically normalised" here
%% means the normalisation that makes ΔΦ = 1 per half-period, fixed by the
%% boundary configuration; M2-7 (bracket convention) remains with the PI.

With these two facts \eqref{eq:2D} is the single-modulus system
\begin{equation}\label{eq:2D-single}
  S_{2}\;=\;\int_{\Ttwo/\Ztwo}\!\sqrt{g_{2}}\,
  \Bigl[\frac{1}{2\kappa_{2}^{2}}R_{2}\;-\;\tfrac12\lvert F\rvert^{2}\;-\;V(a)\Bigr]\,,
\end{equation}
with $V(a)$ the bulk potential assembled from $\Lambda_{12}$, the fibre
curvature and the dilaton potential, continuous on the base. It is the
restriction of the Master Equation and not an axiom. \rigD

\subsection{Bogomol'nyi completion}\label{sec:bps}

Let $\theta$ be the base coordinate transverse to the wall on which the
flux sits, and consider a static profile $a(\theta)$. The energy
functional of \eqref{eq:2D-single} admits a completion of the square once
$V(a)$ balances the geometric bulk terms,
\begin{equation}\label{eq:square}
  \mathcal{E}\;\propto\;\int\sqrt{g_{2}}\,
  \Bigl(\partial_{\theta}a\;\pm\;\frac{\kappa_{2}^{2}}{2}\,\frac{\lvert F_{\psi\omega}\rvert}{\sqrt{g_{2}}}\Bigr)^{2}
  \;\ge\;0\,,
\end{equation}
saturated when the square vanishes. The balance of $V$ against the bulk
terms is the familiar one of warped five-dimensional models
\cite{RS1999}; here it is not a tuning but the condition that
$\Lambda_{12}$ be the value at which the base system admits a solution
with the given topology, which is how $\Lambda_{12}$ is fixed in this
series (Section~\ref{sec:action}). The negative root, a decaying warp,
gives the first-order equation
\begin{equation}\label{eq:bps}
  \partial_{\theta}a\;=\;-\frac{\kappa_{2}^{2}}{2}\,\frac{\lvert F_{\psi\omega}\rvert}{\sqrt{g_{2}}}\,,
\end{equation}
and everything below is an exact consequence of it. \rigR\ (as a theorem of
\eqref{eq:2D-single}); the balance condition is \rigD.
%% Minimality Register M2-6 (Λ₁₂ by stationarity) is where the "tuning" lives.

\subsection{Flux localisation and quantisation}\label{sec:flux}

Dirac quantisation over the orbifold is $\frac{1}{2\pi}\int F=N$. For $N$
units at the drive-fixed corner $P_{1}=(0,0)$ --- the one corner fixed by
both the fold and the drive (M1 Propositions~4.3, 5.6) --- the density is
\begin{equation}\label{eq:flux-density}
  F_{\psi\omega}\;=\;N\pi\,\frac{\delta(\theta_{\psi})\,\delta(\theta_{\omega})}{\sqrt{g_{2}}}\,,
\end{equation}
checked directly: over the fundamental domain
$\frac{1}{2\pi}\int F=N/2$, and the domain is half of the covering torus,
so the total over the cover is $2\cdot N/2=N$ (\texttt{L5\_FD\_check}).
\rigR

\subsection{Solving the first-order equation}\label{sec:solve}

Away from the corner the localised flux vanishes and the flow is driven by
the constant bulk potential, giving the linear profile
\begin{equation}\label{eq:bulk}
  \partial_{\theta}a=-k\qquad\Longrightarrow\qquad a(\theta)=a(0)-k\lvert\theta\rvert\,;
\end{equation}
that the slope is constant between the walls is the boundary condition of
the construction, and the fibre's smooth backreaction, absorbed in $V$,
alters the shape between the walls without touching what follows. At the
corner, integrating \eqref{eq:bps} across $[-\epsilon,\epsilon]$,
\begin{equation}\label{eq:kink}
  [\partial_{\theta}a]_{0}\;=\;-\frac{\kappa_{2}^{2}}{2}\int_{-\epsilon}^{\epsilon}\frac{\lvert F_{\psi\omega}\rvert}{\sqrt{g_{2}}}\,\mathrm{d}\theta
  \;=\;-\kappa_{2}^{2}N\pi\,,
\end{equation}
exactly: each unit of flux contributes a kink of exactly $\kappa_{2}^{2}\pi$,
and no smooth term contributes as $\epsilon\to0$. From \eqref{eq:bulk} the
same kink is $a'(0^{+})-a'(0^{-})=-2k$, so
\begin{equation}\label{eq:balance}
  2k\;=\;\kappa_{2}^{2}\,N\pi\,.
\end{equation}
\rigR\ (the jump); \rigD\ (constant slope).
%% Q16: jump exact, slope a boundary condition. Stated as such.

\subsection{Gauss--Bonnet normalisation}\label{sec:gb}

The half-period $L$ and the two-dimensional coupling are not independent.
By Proposition~\ref{prop:chi-pillowcase} the total curvature $4\pi$ sits as
$\pi$ at each corner; the normalisation connecting the deficit angles to
the length of the domain fixes
\begin{equation}\label{eq:kappa-fix}
  \kappa_{2}^{2}\;=\;\frac{2}{L}\,,
\end{equation}
which, with \eqref{eq:kappa-eff}, is a relation among the scales of the
geometry: $\kappa_{12}^{2}=2V_{4}\mathrm{Vol}_{0}e^{6a_{0}}/L$, the
half-period measured in the gravitational unit weighted by the fibre
volume. It is the one normalisation of this section and we tag it as
such. \rigD
%% Assumption 4 of the WP4 list: κ₂² = 2/L is reproduced as in Paper 2 / the
%% Addendum, with eq:kappa-eff making it a scale relation. A derivation from
%% the orbifold Gauss–Bonnet theorem with the e^{6a} weights would upgrade it
%% to [R]; not supplied here. PI to confirm the reading.

\subsection{The result}\label{sec:result}

Substituting \eqref{eq:kappa-fix} in \eqref{eq:balance},
$2k=(2/L)N\pi$, hence $kL=N\pi$, which is Theorem~\ref{thm:warp-exponent}
(\texttt{L5\_balance}). The smallest case by hand: $N=1$, one unit at one
corner, kink $\kappa_{2}^{2}\pi=2\pi/L$, rate $k=\pi/L$, exponent across
the half-period exactly $\pi$; the general case is $N$ copies of this
balance. Corollary~\ref{cor:zeta} is the composition with
Proposition~\ref{prop:flux}: $\zeta=kL=\chi(\Keight)\pi=12\pi$
(\texttt{L5\_zeta}). Three checks: (i) if the fundamental domain
degenerates to $S^{1}/\Ztwo$ the two-form flux becomes a brane tension
and \eqref{eq:kink} is the Israel junction condition of \cite{RS1999};
(ii) the saturated configuration has positive-definite flux energy;
(iii) $\partial_{\theta}a=-k<0$ between the walls, so the warp decreases
monotonically. \rigR

\subsection{Which direction, and the two sheets}\label{sec:direction}

Nothing above depends on whether the transverse coordinate $\theta$ is
$\theta_{\psi}$ or $\theta_{\omega}$: the normalisation \eqref{eq:kappa-fix}
is applied per direction and \eqref{eq:kL-Npi} is direction-blind. The
sheet exchange $\tau$ of Section~\ref{sec:handedness} fixes $P_{1}$ and
$P_{4}$ and exchanges $P_{2}$ and $P_{3}$, the $\psi$- and
$\omega$-half-periods (\texttt{L5\_tau\_corners}), so a profile transverse
to the wall of one sheet is, under $\tau$, a profile transverse to the wall
of the other. The exponent $12\pi$ is therefore $\tau$-symmetric, and
calling the transverse period $L_{\omega}$ or $L_{\psi}$ is the choice of
which sheet one calls $\psi$, the convention of M1 Appendix~B. What is not
conventional is the aspect ratio: the two half-periods differ by
$\tau_{2}=L_{\omega}/L_{\psi}=1.8466$ (M1 Proposition~8.2), so the rate $k$
differs by $\tau_{2}$ between the sheets, and any quantity built from a
period length rather than from the exponent depends on which sheet is
described. In this series ``$L_{\omega}=12\pi$'' means the half-period
transverse to the $\psi$-wall, in the sheet called $\psi$, in the unit of
\eqref{eq:kappa-fix}. \rigD
%% PI ruling Q15. Register entry 16 (Paper 2 θ_ω vs Addendum θ_ψ).

\subsection{What is derived and what is read}\label{sec:summary4}

The jump \eqref{eq:kink}, the balance \eqref{eq:balance}, the integers of
Section~\ref{sec:integers} and the composition \eqref{eq:zeta} are
theorems (\rigR). Read, and named: the placement of the twelve units on
the base cycle (\rigD, Proposition~\ref{prop:flux}); the constant slope
between walls (\rigD, boundary condition); the normalisation
\eqref{eq:kappa-fix} (\rigD); the tie $a=-\zeta\QCD$ with $\QCD=1$ at the
far wall (\rigD, the definition of the dilaton's normalisation); and the
balance of $V$ as the determination of $\Lambda_{12}$ (\rigD). The value
$\zeta=12\pi$ rests on these and on nothing outside this paper and M1.
